% Pursuing Stacks (cotes 134-2 to 134-8): archive cote 134-4, site page 99 (PDF page 100), typescript page 186.
% typescript. Draft compared: ps3.tex:5040-5102.
% AI-transcribed by Claude Opus 5.5 (claude-opus-5-5, adaptive thinking, effort high), Message Batches API, batch ps2 (msgbatch_01NVwr1Yuo8qJf5rykeGBXPK), 2026-11-12; not yet checked by eye.
\scanpage{134-4/99}
\mspage{186}

If we now look upon the projection map of C$_o$ upon one factor \hand{$\Delta$}$_1$
(carrying a and b  into \hand{$\{$}0\hand{$\}$} and c into \hand{$\{$}1\hand{$\}$}),  we get a functor which is
fibering,  and \xout{xxx} whose fibers are \hand{$\Delta$}$_1$ and \hand{$\Delta$}$_o$ ,  which are $\underline{W}$-aspheric.
\margin{$\times$}Hence \uline{the fibration axiom} L 5)  \uline{on} $\underline{W}$  \uline{implies the Mayer-Vietoris axiom}
L e) (page 171).   This argument rather convinces me that the fibration
axiom should be strong enough to imply all \hand{$\check{\ }$}Cech-type\inserted{W-}asphericity criteria
which one may devise (provided of course they are reasonable,  namely
hold for ordinary weak equivalence !).   More and more,  it seems that
the basic requirements to make upon a \xout{\unsure{modelizer}} basic localizer,  which
will imply maybe all others,  are L 1) (saturation),  L 3')  (the\hand{``}stan-%
dart interval axiom\hand{''} namely  \over{\unsure{4}}{$\Delta$}$_1$  is $\underline{W}$-aspheric ),  and the powerful
fibration axiom L 5) .  This brings to my mind though the condition L a)
of page 165,  namely that  f\hand{$\in$}$\underline{W}$  implies  \hand{$\pi$}$_o$(f)  bijective,  which was'nt
needed really for the famous ``key result'' I was then after,  and which
for this reason I then was looking at almost as something accessory !
I now do feel though that is quite an essential requirement,  even though
we made no formal use of it  (except very incidentally,  in the slightly
weaker form  L a'),  which just means that Hot$_{\hand{$\underline{W}$}}$ is'nt equivalent to the
final category ).   I would therefore add it now to the list of really
basic requirements on a ``basic localizer'',  and rebaptize it therefore
as  L 6) ,  namely:

L 6)  (Connectedness axiom)  f\hand{$\in$}$\underline{W}$  implies  \hand{$\pi$}$_o$(f)  bijective,  i.e.
the functor  \hand{$\pi$}$_o$: (Cat)\hand{$\rightarrow$}(Sets)  factors through  $\underline{W}^{-1}$(Cat)  =
(Hot$_{\underline{W}}$ )  to give rise to a functor
\[ \hand{$\pi$}_o : (\mathrm{Hot}_{\underline{W}}) \hand{$\longrightarrow$} (\mathrm{Sets}) \qquad . \]
This,   as was recalled on page 166,  is more than needed to imply
\[ h_{\underline{W}} = h_{(\mathrm{Cat})} \quad , \]
namely the homotopy structure in (Cat)  associated to $\underline{W}$ \xout{xx} (in terms of
$\underline{W}$-aspheric intervals)  is just the canonical homotopy structure (defined
in terms of 0-connected intervals),  which is also the homotopy structure
defined by the single ``basic'' (multiplicative) interval \hand{$\Delta$}$_1$ .

11.4.

I forgot to clear up still another point of terminology  -  namely
about ``weak homotopy intervals'' -  it turns out finally we never quite
came around defining what ``homotopy intervals'' which are'nt weak should
be !   T\over{\unsure{n}}{h}e situation is very silly indeed  -  so henceforth I'll just

\margin{99}
